% =====================================================================
%  CHAPTER 7: 儿童少年常见病（对应大纲第七章）
%  内容来源：往届笔记第九章，参考PPT第八章
% =====================================================================
\chapter{儿童少年常见病}
\setcounter{chapter}{7}
\chapterintro{
\textbf{掌握：}儿童少年常见病的防治：单纯性肥胖、视力低下、龋齿、脊柱弯曲异常的发病特点或流行规律、防控措施。\newline
\textbf{熟悉：}健康的含义。
}

% ===== 第一部分：熟悉内容 =====
\section{概述}

\subsection{健康的含义}

\begin{defbox}
\textbf{健康}不仅是没有疾病或不虚弱，而是身体、心理和社会适应的完好状态。

\textbf{衡量健康的维度和指标体系：}
\begin{enumerate}[leftmargin=*]
    \item \textbf{生理健康：}体格发育、生理功能、运动能力、日常生活自理能力。
    \item \textbf{心理健康：}心理/精神症状和行为表现、认知功能、智力和个性、主观幸福度。
    \item \textbf{社会健康。}
    \item \textbf{道德健康。}
\end{enumerate}
\end{defbox}

\subsection{儿童少年常见病概述}

\begin{defbox}
\textbf{儿童少年常见病（common diseases in children and adolescents）：}在儿童少年中发病率较高的一类疾病，发病与其正处于生长发育阶段有关，也受幼儿园和学校集体生活影响。
\end{defbox}

\begin{keybox}
\textbf{中国儿童少年常见病流行特点：}

\begin{enumerate}[leftmargin=*]
    \item \textbf{发病率基本得到控制的常见病：}沙眼、蛔虫感染。
    \item \textbf{患病率居高不下的常见病：}近视。
    \item \textbf{检出率持续上升的常见病：}超重肥胖。
    \item \textbf{发展不平衡造成的常见病：}营养不良。
\end{enumerate}

\textbf{危险因素：}

\begin{enumerate}[leftmargin=*]
    \item \textbf{遗传因素：}单基因遗传（高度近视）；多基因遗传（超重肥胖、单纯性近视，近视遗传度为65-70\%）。
    \item \textbf{环境因素：}环境内分泌干扰物（EDCs），家庭环境（生活环境、家庭结构、家庭经济状况、父母文化程度、教育方式）。
    \item \textbf{个体因素：}膳食摄入、饮食习惯、体力活动。
\end{enumerate}

\textbf{预防控制总策略：}

\begin{enumerate}[leftmargin=*]
    \item \textbf{健康促进：}制定并认真落实学校健康促进政策；保证学生每天体育锻炼、户外活动时间；定期组织体检，建立学生健康档案。
    \item \textbf{环境改善：}改善教室环境；创造良好家庭和学习环境。
    \item \textbf{健康素养提升：}提高对常见病认知程度和防范意识；养成良好行为和生活方式。
\end{enumerate}
\end{keybox}

% ===== 第二部分：掌握内容 =====
\section{儿童少年肥胖}\difficulty

\begin{defbox}
\textbf{肥胖（obesity）：}在遗传和环境因素交互作用下，因能量摄入>能量消耗，导致体内脂肪积聚过多，危害健康的慢性代谢性疾病。

WHO标准：超重 25≤BMI≤30，肥胖 BMI≥30。
中国标准：超重 24≤BMI≤28，肥胖 BMI≥28。
\end{defbox}

\subsection{危险因素}

\begin{defbox}
\textbf{继发性肥胖：}病因明确，与原发性病因有关，需先治疗其原发性病因。

\textbf{单纯性肥胖：}在遗传和环境因素交互作用下，因能量摄入>能量消耗，导致体内脂肪积聚过多，危害健康的慢性代谢性疾病。肥胖是多基因关联累加结果，属于多基因遗传。
\end{defbox}

\begin{keybox}
\textbf{单纯性肥胖的敏感期：}孕后期、婴儿期、青春早期、青春后期（无幼儿期和青春中期）。

\textbf{1. 遗传因素}
\begin{enumerate}[leftmargin=*]
    \item 是重要因素，但不是决定因素，肥胖是多基因关联累加结果，属于多基因遗传。
    \item 瘦素通过3种途径调节机体脂肪沉积：①抑制食欲；②增加能量消耗；③抑制脂肪合成。
\end{enumerate}

\textbf{2. 个人因素}
\begin{enumerate}[leftmargin=*]
    \item \textbf{膳食摄入：}膳食热能过高；膳食结构不合理；食用能量密度高的食物。
    \item \textbf{饮食行为：}不良食品选择倾向（甜食、甜饮料等）；不良饮食方式（煎、炸烹饪方式等）；不良膳食制度（不吃早餐等）；不良食物奖惩形式（吃西式快餐作为奖励）。
    \item \textbf{体力活动：}体育活动减少；静态生活时间增加。
    \item \textbf{肠道菌群。}
\end{enumerate}

\textbf{3. 环境因素}
\begin{enumerate}[leftmargin=*]
    \item \textbf{社会环境因素：}传统健康观念"胖是健康，以胖为福"；肥胖易感环境（obesogenic environment）：全球化、城市化带来的膳食快餐化、强大市场诱惑、媒体对饮食行为的影响。
    \item \textbf{家庭环境因素：}受母亲文化程度等因素影响，文化程度高的母亲更倾向于调整孩子的饮食量，鼓励子女参加体育活动。
\end{enumerate}
\end{keybox}

\subsection{肥胖危害}

\begin{keybox}
\textbf{1. 健康危害：}心血管疾病、早期动脉粥样硬化、2型糖尿病、代谢综合征、非酒精性脂肪肝、阻塞性睡眠呼吸暂停（OSA）。

\textbf{2. 心理危害和行为危害}
\begin{enumerate}[leftmargin=*]
    \item 心理危害：心理压抑，心理精神障碍（如抑郁和自杀意念等）。
    \item 行为危害：饮食失调症、社会交往能力下降、吸烟饮酒行为、校园欺负行为，甚至自杀意念和行为。
\end{enumerate}

\textbf{3. 社会危害}
\begin{enumerate}[leftmargin=*]
    \item 直接成本：控制和治疗肥胖付出的成本。
    \item 机会性成本：肥胖相关疾病或过早死亡造成的医疗成本或经济损失。
    \item 间接成本：个人和社区的间接（社会）负担，如病假、个人用于减轻体重的花费等。
\end{enumerate}
\end{keybox}

\subsection{肥胖筛查}

\begin{keybox}
\textbf{筛查方法：}
\begin{enumerate}[leftmargin=*]
    \item 目测法
    \item 身高标准体重法：WHO推荐使用，限于0-6岁儿童
    \item 体质量指数（BMI）：简便易测，能进行跨人群（现状）和跨时段（趋势）分析
    \item 腹部脂肪测量：对预测肥胖造成的危害和儿科临床应用有重要意义和价值
\end{enumerate}

\textbf{儿童少年肥胖筛查标准：}

\begin{table}[H]
\centering
\caption{儿童少年肥胖筛查标准比较}
\begin{tabularx}{\textwidth}{l>{\centering\arraybackslash}p{2.5cm}>{\centering\arraybackslash}p{2.5cm}>{\centering\arraybackslash}p{2.5cm}>{\centering\arraybackslash}p{2.5cm}}
\hline
& \textbf{NCHS标准} & \textbf{IOTF标准} & \textbf{中国标准} \\
\hline
超重 & P85≤BMI≤P95 & 25≤BMI≤30 & 24≤BMI≤28 \\
\hline
肥胖 & BMI≥P95 & BMI≥30 & BMI≥28 \\
\hline
\end{tabularx}
\end{table}

\textbf{体脂率（body fat percentage, \%BF）：}人体脂肪组织占体重百分比，是较直观判断肥胖的指标。
\end{keybox}

\subsection{预防控制}

\begin{keybox}
\textbf{1. 普遍性预防：}运用健康教育、健康促进理论，依托建设健康学校、提高学生健康素养的形式，通过制定政策、创建支持环境、社区积极参与、开展健康教育、提供健康咨询和指导，培养儿童少年健康行为习惯和生活方式。

\textbf{2. 针对性预防}
\begin{enumerate}[leftmargin=*]
    \item 平衡膳食、建立良好膳食制度、保持食物多样化、三餐热量合理分配。
    \item 培养健康饮食行为，少吃油炸食品、不喝含糖饮料等。
    \item 积极参加体育活动（每天60min以上中高等强度体育活动），减少静态活动时间（每天看电视玩游戏≤2h）。
    \item 不提倡节食、挑食，不使用减肥药、腹泻等不健康"减肥"方式。
\end{enumerate}

\textbf{3. 干预性防控}
\begin{enumerate}[leftmargin=*]
    \item \textbf{饮食调整：}膳食结构合理，减少脂肪摄入。
    \item \textbf{身体活动指导：}坚持规律有氧运动，通过评估调整运动方案。
    \item \textbf{行为矫正：}制定行为改变目标，以日记等方式进行自我监督，正向鼓励良好行为，家长以身示范。
    \item \textbf{心理疏导：}树立正常健康观和意识，疏解压抑、自卑、消极心理。
\end{enumerate}
\end{keybox}

\section{儿童少年近视}\difficulty

\begin{defbox}
\textbf{近视：}不使用调节功能状态下，远处来的平行光线在视网膜感光层前方聚焦。

诊断标准：（1）近视：等效球镜度数≤-0.50D；（2）高度近视：等效球镜度数≤-6.0D（600度及以上）。

分类：（1）按屈光度：低度（-0.25D~-3.0D）、中度（-3.25D~-6.0D）、高度（-6.25D~-9.0D）；（2）按有无调节因素参与：假性、真性、半真性；（3）按屈光要素改变：轴性、屈光性。
\end{defbox}

\subsection{发生机制}

\begin{keybox}
\textbf{正视化（emmetropization）：}大部分儿童出生时远视，随发育进程眼轴变长、晶状体和角膜弯曲度变平、眼睛屈光度变强，发展为正视的过程。

\textbf{近视发生机制：}学习时间过长，光照条件不良→眼睛常处于高度调节紧张状态→看远处时睫状肌仍处于收缩状态→悬韧带放松→晶状体屈光力过强→近视。

\begin{enumerate}[leftmargin=*]
    \item \textbf{假性近视（pseudo myopia）：}因过度视近工作形成的近视为调节紧张性近视，属功能性改变，是近视防控关键窗口期。采取视力保护措施→纠正不良用眼习惯，放松睫状肌→视力可恢复正常。
    \item \textbf{真性/轴性近视：}假性近视后依然不注意用眼卫生→引起眼球充血、眼压升高、眼球壁弹性下降→刺激眼轴伸长。
\end{enumerate}

\textbf{生理基础：}正视化完成后角膜曲率基本不再变化，但眼轴长度可继续生长。
\end{keybox}

\subsection{影响因素与预防控制}

\begin{keybox}
\textbf{影响因素：}1.遗传因素；2.环境因素；3.体质健康因素。

\textbf{预防控制：}

\textbf{1. 近视预防}
\begin{enumerate}[leftmargin=*]
    \item 加强学校预防近视工作
    \item 培养良好用眼习惯
    \item 创造良好生活环境
    \item 认真做好眼保健操
    \item 积极参加户外活动，亲近阳光
\end{enumerate}

\textbf{2. 近视矫正}
\begin{enumerate}[leftmargin=*]
    \item 及时检查
    \item 合理配镜
    \item 抗胆碱能药物
    \item 手术治疗
\end{enumerate}
\end{keybox}

\begin{exambox}
△ 近视每年检查2次。
\end{exambox}

\section{龋齿和牙周病}\difficulty

\begin{defbox}
\textbf{龋齿（dental caries）：}牙齿在身体内外因素作用下，硬组织脱矿，有机质溶解，牙组织进行性破坏，导致牙齿缺损的儿童常见病。

△流行状况指标：龋患率；严重程度：龋均、患者龋均；治疗状况：DMFT。
\end{defbox}

\subsection{发生原因与危害}

\begin{keybox}
\textbf{发生原因（四联致病因素）：}
\begin{enumerate}[leftmargin=*]
    \item \imp{细菌和菌斑}
    \item \imp{宿主因素}
    \item \imp{食物因素}
    \item \imp{时间因素}
\end{enumerate}

\textbf{危害：}
\begin{enumerate}[leftmargin=*]
    \item \textbf{局部：}影响食欲，干扰咀嚼、消化和吸收，导致营养缺乏，且随龋病发展可引发牙髓炎、颜面蜂窝织炎、根周脓肿、齿槽溢脓等严重口腔疾病。
    \item \textbf{全身性疾病：}通过变态反应引发风湿性关节炎、心内膜炎、肾炎。
    \item 损害个人形象，影响心理健康。
\end{enumerate}
\end{keybox}

\subsection{牙周病}

\begin{tipbox}
\textbf{分类：}
\begin{enumerate}[leftmargin=*]
    \item 牙龈病：牙龈组织
    \item 牙周炎：牙龈组织+深层组织
\end{enumerate}

\textbf{发生原因：}细菌，宿主因素，环境因素。
\end{tipbox}

\subsection{预防控制}

\begin{keybox}
\textbf{预防控制措施：}
\begin{enumerate}[leftmargin=*]
    \item \textbf{加强口腔保健：}定期检查，早期诊断
    \item \textbf{控制牙菌斑：}刷牙3min
    \item \textbf{注意饮食卫生，增强宿主抗龋力}
    \item \textbf{加强健康教育}
    \item \textbf{健全学校口腔疾病防治网}
\end{enumerate}
\end{keybox}

\begin{exambox}
△ 龋齿每年检查1次。
\end{exambox}

\section{脊柱弯曲异常}\difficulty

\begin{defbox}
\textbf{脊柱弯曲异常（vertebral column defects）：}脊柱弯曲超出正常生理范围。
\end{defbox}

\begin{keybox}
\textbf{发生原因：}
\begin{enumerate}[leftmargin=*]
    \item \textbf{习惯性姿势：}不良站姿，不良坐姿，歪头写字
    \item \textbf{桌椅高矮不适合}
    \item \textbf{体力活动缺乏}
    \item \textbf{营养和体质因素}
\end{enumerate}

\textbf{预防控制：}
\begin{enumerate}[leftmargin=*]
    \item 加强健康教育
    \item 保证足够体育锻炼和体力活动
    \item 调整座椅高度
    \item 定期筛查，早期发现
    \item 脊柱弯曲异常矫正
\end{enumerate}
\end{keybox}

% ----- 复习题 -----
\reviewcontent{
\section{复习题}

\begin{enumerate}[leftmargin=*, itemsep=8pt]
    \item \textbf{【名词解释】}儿童少年常见病；肥胖；继发性肥胖；单纯性肥胖；近视；正视化；假性近视；龋齿；脊柱弯曲异常；体脂率
    \item \textbf{【简答题】}简述中国儿童少年常见病的流行特点。
    \item \textbf{【简答题】}简述单纯性肥胖的危险因素和敏感期。
    \item \textbf{【简答题】}简述肥胖的三级预防控制策略。
    \item \textbf{【简答题】}简述近视的发生机制及假性近视与真性近视的区别。
    \item \textbf{【简答题】}简述龋齿的四联致病因素和预防控制措施。
    \item \textbf{【简答题】}简述脊柱弯曲异常的发生原因和预防措施。
\end{enumerate}

\subsection*{参考答案要点}

\begin{enumerate}[leftmargin=*, itemsep=5pt]
    \item \textbf{儿童少年常见病}：在儿童少年中发病率较高的一类疾病，发病与生长发育阶段和集体生活影响有关。\textbf{肥胖}：在遗传和环境因素交互作用下，因能量摄入>能量消耗，导致体内脂肪积聚过多，危害健康的慢性代谢性疾病。\textbf{继发性肥胖}：病因明确，与原发性病因有关，需先治疗其原发性病因。\textbf{单纯性肥胖}：在遗传和环境因素交互作用下，因能量摄入>能量消耗，导致体内脂肪积聚过多，属于多基因遗传。\textbf{近视}：不使用调节功能状态下，远处来的平行光线在视网膜感光层前方聚焦。\textbf{正视化}：大部分儿童出生时远视，随发育进程发展为正视的过程。\textbf{假性近视}：因过度视近工作形成的调节紧张性近视，属功能性改变，是近视防控关键窗口期。\textbf{龋齿}：牙齿硬组织脱矿，有机质溶解，牙组织进行性破坏导致牙齿缺损的儿童常见病。\textbf{脊柱弯曲异常}：脊柱弯曲超出正常生理范围。\textbf{体脂率}：人体脂肪组织占体重百分比，是较直观判断肥胖的指标。

    \item (1)发病率基本得到控制：沙眼、蛔虫感染；(2)患病率居高不下：近视；(3)检出率持续上升：超重肥胖；(4)发展不平衡造成：营养不良。

    \item 敏感期：孕后期、婴儿期、青春早期、青春后期（无幼儿期和青春中期）。危险因素：(1)遗传（多基因遗传，瘦素通过抑制食欲/增加能耗/抑制脂肪合成调节）；(2)个人因素（膳食摄入、饮食行为、体力活动不足、肠道菌群）；(3)环境因素（社会环境——传统观念和肥胖易感环境；家庭环境——母亲文化程度等）。

    \item (1)普遍性预防：运用健康教育和健康促进理论，培养健康行为习惯和生活方式；(2)针对性预防：平衡膳食、培养健康饮食行为、积极参加体育活动（每天60min以上）、减少静态活动时间（≤2h/天）；(3)干预性防控：饮食调整、身体活动指导、行为矫正、心理疏导。

    \item 机制：学习时间过长+光照不良→眼睛高度调节紧张→看远时睫状肌仍收缩→悬韧带放松→晶状体屈光力过强→近视。假性近视：调节紧张性近视，功能性改变，可恢复正常（近视防控关键窗口期）。真性近视：假性近视后不注意用眼卫生→眼球充血、眼压升高、眼球壁弹性下降→眼轴伸长。

    \item 四联致病因素：细菌和菌斑、宿主因素、食物因素、时间因素。预防：(1)加强口腔保健，定期检查，早期诊断；(2)控制牙菌斑，刷牙3min；(3)注意饮食卫生，增强宿主抗龋力；(4)加强健康教育；(5)健全学校口腔疾病防治网。

    \item 原因：(1)习惯性姿势（不良站姿、坐姿、歪头写字）；(2)桌椅高矮不适合；(3)体力活动缺乏；(4)营养和体质因素。预防：(1)加强健康教育；(2)保证足够体育锻炼；(3)调整座椅高度；(4)定期筛查，早期发现；(5)脊柱弯曲异常矫正。
\end{enumerate}
}

\newpage
